\chapter{Słabsze zasady wyboru i ich ograniczenia}\label{ch:weak-choice}
\chaptermark{Słabsze zasady wyboru}
\section{Dlaczego rozróżniamy postacie wyboru?}
Aksjomat wyboru (AC) dotyczy dowolnej rodziny niepustych zbiorów. W wielu dowodach potrzebujemy jednak wyboru tylko z rodziny przeliczalnej albo kolejnych wyborów zależnych od poprzednich. W tym rozdziale podstawową teorią jest ZF; wszystkie dodatkowe założenia podajemy jawnie.
\begin{definicja}[Wybór przeliczalny]\label{def:ac-omega}
Zasada $\mathrm{AC}_{\omega}$ głosi, że dla każdej funkcji $n\mapsto A_n$ o dziedzinie $\omega$, której wartości są niepustymi zbiorami, istnieje funkcja $f$ o dziedzinie $\omega$ taka, że $f(n)\in A_n$ dla każdego $n\in\omega$.
\end{definicja}
\begin{definicja}[Wybór zależny]\label{def:dc}
Zasada wyboru zależnego $\mathrm{DC}$ głosi: dla każdego niepustego zbioru $X$, każdej relacji $R\subseteq X\times X$ spełniającej $\forall x\in X\ \exists y\in X\ (xRy)$ i każdego $x_0\in X$ istnieje funkcja $f:\omega\to X$ taka, że $f(0)=x_0$ i $f(n)Rf(n+1)$ dla wszystkich $n\in\omega$.
\end{definicja}
\begin{twierdzenie}[Łańcuch implikacji]\label{thm:choice-chain}
W ZF zachodzą implikacje $\mathrm{AC}\Rightarrow\mathrm{DC}\Rightarrow\mathrm{AC}_{\omega}$.
\end{twierdzenie}
\begin{proof}
Załóżmy AC. Dla każdego $x\in X$ niepusty jest zbiór $R[x]=\{y\in X:xRy\}$. AC daje funkcję $c:X\to X$ z $c(x)\in R[x]$. Zwykła rekurencja na $\omega$ z $f(0)=x_0$ i $f(n+1)=c(f(n))$ daje DC.

Załóżmy DC i rodzinę $(A_n)_{n\in\omega}$ niepustych zbiorów. Niech $X$ będzie zbiorem wszystkich skończonych funkcji $s$ o dziedzinie pewnego $n\in\omega$ spełniających $s(k)\in A_k$ dla $k<n$. Zbiór $X$ istnieje w ZF: wszystkie takie funkcje są podzbiorami $\omega\times\bigcup_{n\in\omega}A_n$, więc $X$ otrzymujemy przez wyróżnianie w zbiorze potęgowym. Nie zakładamy istnienia elementu dowolnego nieskończonego produktu. Jest niepusty, bo zawiera funkcję pustą. Niech $sRt$ oznacza, że $t$ przedłuża $s$ o dokładnie jedną parę $(n,a)$ z $a\in A_n$, gdzie $n=\dom s$. Każdy $s$ ma następnika, ponieważ $A_n\ne\varnothing$. Z DC istnieje ciąg $s_0,s_1,\ldots$ zaczynający się od funkcji pustej i spełniający $s_nRs_{n+1}$. Wówczas $f=\bigcup_{n\in\omega}s_n$ jest funkcją wyboru dla całej rodziny.
\end{proof}
\begin{uwaga}[Nie mylić dwóch rodzajów wyboru]
Z faktu, że każdy $A_n$ jest niepusty, wynika dla każdego ustalonego $n$ istnienie elementu $a_n$. Nie wynika z samej logiki pierwszego rzędu istnienie jednej funkcji $n\mapsto a_n$. Właśnie tę lukę wypełnia $\mathrm{AC}_{\omega}$.
\end{uwaga}
\section{Zastosowanie: przeliczalna suma zbiorów przeliczalnych}
\begin{twierdzenie}\label{thm:countable-union-cc}
W ZF+$\mathrm{AC}_{\omega}$ dla każdej rodziny $(A_n)_{n\in\omega}$ zbiorów co najwyżej przeliczalnych suma $\bigcup_n A_n$ jest co najwyżej przeliczalna.
\end{twierdzenie}
\begin{proof}
Dla każdego $n$ niech $E_n$ będzie zbiorem wszystkich iniekcji $A_n\to\omega$. Jest on niepusty z założenia. Z $\mathrm{AC}_{\omega}$ wybieramy $e_n\in E_n$ jednocześnie dla wszystkich $n$. Każdemu $a\in\bigcup_n A_n$ przypiszmy parę $(n,e_n(a))$, gdzie $n$ jest najmniejszym indeksem z $a\in A_n$. Jest to iniekcja do $\omega\times\omega$, który jest przeliczalny na mocy Twierdzenia~\ref{thm:card-n2}. Nie potrzeba wybierać elementów z pustych zbiorów.
\end{proof}
\section{Słabe zasady a zwykłe konstrukcje}
\begin{przyklad}[Skończenie wiele wyborów]
Jeżeli $A$ i $B$ są niepuste, to istnieją $a\in A$ i $b\in B$, a więc para $(a,b)\in A\times B$. Ten argument wykorzystuje dwukrotnie regułę wprowadzania kwantyfikatora egzystencjalnego, nie AC. Dla nieskończenie wielu zbiorów analogiczne rozumowanie nie tworzy automatycznie funkcji wyboru.
\end{przyklad}
\begin{uwaga}[Granice twierdzeń]
Implikacje odwrotne do łańcucha z Twierdzenia~\ref{thm:choice-chain} nie są twierdzeniami ZF (przy standardowych założeniach niesprzeczności). Wykazanie tego wymaga budowania modeli teorii mnogości; nie traktujemy samego braku znanego dowodu jako dowodu niezależności. Metody modelowe wprowadzimy w dalszych rozdziałach.
\end{uwaga}

\section{Różnica między wyborem a rekurencją}
\begin{przyklad}[Jawna reguła nie wymaga wyboru]
Niech $A_n=\{n,n+1\}$ dla $n\in\omega$. Funkcja $f(n)=n$ jest funkcją wyboru już w ZF: wskazaliśmy jeden wzór określający wszystkie wartości. Zasada $\mathrm{AC}_\omega$ staje się istotna dopiero wtedy, gdy nie ma z góry zadanej reguły wyboru.
\end{przyklad}
\begin{uwaga}[Wybór dla rodzin indeksowanych]
Zasada $\mathrm{AC}_\omega$ nie mówi, że dowolna rodzina zbiorów przeliczalnych ma funkcję wyboru; wymaga, aby indeksowanie rodziny było funkcją o dziedzinie $\omega$. Samo stwierdzenie, że każdy zbiór rodziny jest przeliczalny, nie dostarcza indeksowania całej rodziny.
\end{uwaga}


\section{Wybór zależny i twierdzenie Baire'a dla przestrzeni zupełnych}
\label{sec:weak-choice-baire}
W przestrzeni Baire'a $\omega^\omega$ dopuszczalne
przedłużenia skończonych ciągów można wybierać
według ustalonego porządku liczb naturalnych.
W ogólnej zupełnej przestrzeni metrycznej nie ma jednak
kanonicznego sposobu wskazania środka następnej kuli.
Poniższy dowód pokazuje, gdzie wystarcza $\mathrm{DC}$.

\begin{twierdzenie}[Twierdzenie Baire'a z DC]
\label{thm:weak-choice-baire-dc}
W teorii ZF+$\mathrm{DC}$ niech $(X,d)$ będzie
niepustą zupełną przestrzenią metryczną.
Jeśli $(D_n)_{n\in\omega}$ jest ciągiem
zbiorów otwartych i gęstych w $X$,
to $\bigcap_{n\in\omega}D_n$ jest gęsty w $X$.
W szczególności żaden niepusty zbiór otwarty
w $X$ nie jest pierwszej kategorii w $X$.
\end{twierdzenie}
\begin{proof}
\emph{Krok 1: zbiory przybliżeń.}
Ustalmy niepusty zbiór otwarty $U\subseteq X$.
Piszemy
$B(z,r):=\{x\in X:d(x,z)<r\}$
oraz
$\overline B(z,r):=\{x\in X:d(x,z)\le r\}$.
Przez $\overline B(z,r)$ rozumiemy tu
\emph{kulę domkniętą}, która w ogólnej
przestrzeni metrycznej nie musi być
domknięciem topologicznym $B(z,r)$.
Dodatnie promienie wybieramy ze zbioru
$\mathbb Q_{>0}$.

Niech $P$ będzie zbiorem skończonych ciągów
$s=((z_0,r_0),\ldots,(z_{m-1},r_{m-1}))$
elementów $X\times\mathbb Q_{>0}$ spełniających
$r_i<2^{-(i+2)}$ oraz
\[
\overline B(z_0,r_0)\subseteq U\cap D_0,\qquad
\overline B(z_{i+1},r_{i+1})
 \subseteq B(z_i,r_i)\cap D_{i+1}
\quad(i+1<m).
\]
Dopuszczamy także ciąg pusty.
Wszystkie skończone ciągi par są elementami
zbioru
$\bigcup_{m\in\omega}(X\times\mathbb Q_{>0})^m$,
istniejącego w ZF przez rekursję,
zastępowanie i sumę.
Zatem $P$ istnieje przez separację
i zawiera ciąg pusty.

\emph{Krok 2: każde przybliżenie można przedłużyć.}
Jeśli $s$ jest pusty,
otwarty zbiór $U\cap D_0$ jest niepusty
(gdyż $D_0$ jest gęsty).
Wybieramy jeden punkt $z_0$ tego zbioru.
Otwartość daje $\varepsilon>0$ z
$B(z_0,\varepsilon)\subseteq U\cap D_0$.
Wybieramy jedną wymierną liczbę
$0<r_0<\min(\varepsilon,1/4)$;
wtedy
$\overline B(z_0,r_0)\subseteq U\cap D_0$.

Jeśli $s$ ma długość $m>0$, zbiór
$B(z_{m-1},r_{m-1})\cap D_m$
jest niepusty i otwarty.
Jego niepustość wynika z gęstości $D_m$
i niepustości kuli otwartej.
Wybieramy jeden punkt $z_m$ i dodatni
promień wymierny
$r_m<2^{-(m+2)}$, tak mały, aby
\[
\overline B(z_m,r_m)
\subseteq B(z_{m-1},r_{m-1})\cap D_m.
\]
Taki promień istnieje dzięki otwartości
zbioru po prawej i gęstości $\mathbb Q$ w $\mathbb R$.
Każdy z powyższych pojedynczych wyborów
jest zwykłym użyciem kwantyfikatora egzystencjalnego.

\emph{Krok 3: zastosowanie DC.}
Dla $s,t\in P$ niech $sRt$ oznacza, że
$t$ powstaje z $s$ przez dołączenie dokładnie
jednej pary. Krok 2 daje
$\forall s\in P\;\exists t\in P\;(sRt)$.
Zasada $\mathrm{DC}$ zastosowana do $P,R$
i ciągu pustego daje funkcję
$k\mapsto s_k$ o dziedzinie $\omega$,
zaczynającą się od ciągu pustego,
taką że $s_kRs_{k+1}$.
Suma tych zgodnych przybliżeń wyznacza
jednoznaczny ciąg par
$(z_n,r_n)_{n\in\omega}$.

Dla $j>i$ z zagnieżdżenia kul mamy
$z_j\in B(z_i,r_i)$, a więc
$d(z_i,z_j)<r_i<2^{-(i+2)}$.
Zatem $(z_n)$ jest ciągiem Cauchy'ego.
Zupełność daje jego granicę $x\in X$.

\emph{Krok 4: przynależność granicy.}
Dla każdego $i$ i wszystkich $j>i$ zachodzi
$z_j\in\overline B(z_i,r_i)$.
Ponieważ ta kula jest domknięta,
$x\in\overline B(z_i,r_i)$.
Dla $i=0$ mamy zatem $x\in U\cap D_0$,
a dla $i>0$:
\[
x\in\overline B(z_i,r_i)
 \subseteq B(z_{i-1},r_{i-1})\cap D_i.
\]
Wynika stąd
$x\in U\cap\bigcap_{i\in\omega}D_i$.
Ponieważ niepusty otwarty $U$
był dowolny, przekrój zbiorów $D_i$
jest gęsty.

\emph{Krok 5: kategoria pierwsza.}
Jeśli $U$ byłby zawarty w przeliczalnej
sumie zbiorów nigdziegęstych $(F_n)$,
to $D_n:=X\setminus\overline{F_n}$
byłyby otwarte i gęste.
Z poprzedniej części otrzymalibyśmy punkt
\[
x\in U\cap\bigcap_nD_n
\subseteq U\setminus\bigcup_nF_n,
\]
sprzecznie z $U\subseteq\bigcup_nF_n$.
\end{proof}

\begin{przyklad}[Dlaczego wymagana jest zupełność?]
\label{ex:weak-choice-baire-noncomplete}
Przestrzeń $(\mathbb Q,|\cdot|)$ nie jest zupełna:
ciąg wymiernych przybliżeń do $\sqrt2$ jest
ciągiem Cauchy'ego bez granicy w $\mathbb Q$.
Ustalmy bijekcję $e:\omega\to\mathbb Q$,
której istnienie w ZF wynika z konstrukcji
w Twierdzeniu~\ref{thm:card-n2}
i z jawnego kodowania liczb wymiernych
w dowodzie Twierdzenia~\ref{thm:card-continuum}.
Połóżmy
$D_n:=\mathbb Q\setminus
\{e(0),\ldots,e(n)\}$.
Każdy $D_n$ jest otwarty w $\mathbb Q$,
gdyż usuwamy skończenie wiele punktów.
Jest też gęsty, ponieważ
każdy niepusty przedział wymierny zawiera
nieskończenie wiele liczb wymiernych.
Jednak $\bigcap_nD_n=\varnothing$,
ponieważ każdy element $\mathbb Q$
należy do obrazu $e$.
Zatem bez zupełności twierdzenie może zawieść.
\end{przyklad}

\begin{uwaga}[Dokładna rola zasad wyboru]
\label{rem:weak-choice-baire-dc}
W dowodzie Twierdzenia~\ref{thm:weak-choice-baire-dc}
kolejna dopuszczalna kula zależała od poprzedniej.
Jest to właśnie sytuacja opisywana przez DC,
a nie niezależny wybór z gotowej rodziny.
W przestrzeni $\omega^\omega$
dla odpowiednio kodowanych przybliżeń
można wybierać najmniejsze dopuszczalne
przedłużenia, co daje osobny argument w ZF.
Nie wyprowadzamy stąd nieudowodnionej
równoważności DC z twierdzeniem Baire'a
dla wszystkich klas przestrzeni.
\end{uwaga}

\section{Algebry Boole'a, filtry i twierdzenie o ideale pierwszym}
\label{sec:weak-choice-bpi}
W aksjomacie wyboru możemy wskazywać elementy z dowolnej rodziny
niepustych zbiorów. Istnieje jednak słabsza zasada wystarczająca
do wielu zastosowań topologicznych: pozwala \emph{przedłużyć już
spójny układ decyzji} dotyczących zbiorów do decyzji całkowitej.
Zaczniemy od wyjaśnienia wszystkich występujących tu pojęć.

\begin{definicja}[Algebra Boole'a i filtr]\label{def:bpi-boolean-filter}
\emph{Algebrą Boole'a} nazywamy niepusty
zbiór $B$ z dwuargumentowymi działaniami
$\wedge,\vee:B^2\to B$,
jednoargumentowym $\neg:B\to B$
oraz wyróżnionymi elementami $0_B,1_B\in B$.
Dla wszystkich $a,b,c\in B$ zakładamy
łączność i przemienność obu działań,
prawa pochłaniania
\[
a\wedge(a\vee b)=a,\qquad
a\vee(a\wedge b)=a,
\]
oba prawa rozdzielności
\[
a\wedge(b\vee c)=(a\wedge b)\vee(a\wedge c),
\quad
a\vee(b\wedge c)=(a\vee b)\wedge(a\vee c),
\]
oraz $a\wedge\neg a=0_B$
i $a\vee\neg a=1_B$.
W szczególności wynikają stąd
$a\wedge1_B=a$, $a\vee0_B=a$
i idempotentność obu działań.
Definiujemy $a\leq b$ wtedy i tylko wtedy, gdy $a\wedge b=a$.
W całej sekcji zakładamy $0_B\ne1_B$.
\emph{Filtr właściwy} $F\subseteq B$ spełnia:
$1_B\in F$, $0_B\notin F$,
$b,c\in F\implies b\wedge c\in F$ oraz
$b\in F,\ b\leq c\implies c\in F$.
Filtr $U$ jest \emph{ultrafiltrem}, jeżeli jest właściwy i
maksymalny wśród właściwych filtrów względem zawierania.
\emph{Ideał pierwszy} to zbiór $I\subseteq B$, którego
dopełnienie $B\setminus I$ jest ultrafiltrem.
\end{definicja}
\begin{przyklad}[Co oznaczają operacje?]
Dla dowolnego niepustego zbioru $S$
algebra $\mathcal P(S)$ ma
$X\wedge Y=X\cap Y$, $X\vee Y=X\cup Y$,
$\neg X=S\setminus X$, $0_B=\varnothing$, $1_B=S$.
Dla ustalonego $s\in S$ rodzina
$U_s=\{X\subseteq S:s\in X\}$ jest ultrafiltrem.
Filtr zbiorów współskończonych w $\N$ zawiera każdy zbiór
o skończonym dopełnieniu; jest właściwy, ale nie jest
ultrafiltrem, gdyż ani liczby parzyste, ani ich dopełnienie
nie są współskończone.
\end{przyklad}

\begin{lemat}[Kryterium ultrafiltru]\label{lem:bpi-ultrafilter-criterion}
Dla każdego właściwego filtru $U$ na algebrze Boole'a $B$
następujące warunki są równoważne:
(i) $U$ jest ultrafiltrem;
(ii) dla każdego $b\in B$ dokładnie jeden z elementów
$b,\neg b$ należy do $U$.
\end{lemat}
\begin{proof}
Najpierw oba elementy nie mogą należeć do $U$, gdyż
ich iloczyn to $0_B$. Załóżmy, że $U$ jest maksymalny i $b\notin U$.
Filtr generowany przez $U\cup\{b\}$ byłby
właściwym filtrem ściśle większym od $U$,
gdyby wszystkie przekroje $u\wedge b$ ($u\in U$)
były niezerowe. Maksymalność $U$ wyklucza tę
możliwość, więc istnieje $u\in U$ z
$u\wedge b=0_B$. Wtedy $u\leq\neg b$,
a domknięcie filtru ku górze daje $\neg b\in U$.
Dokładnie jeden spośród $b,\neg b$ należy do $U$.
Dla (ii)$\Rightarrow$(i) niech $V\supsetneq U$ będzie
filtrem i wybierzmy $b\in V\setminus U$. Wtedy
$\neg b\in U\subseteq V$, więc $0_B=b\wedge\neg b\in V$.
Zatem $V$ nie jest właściwy i $U$ jest maksymalny.
\end{proof}

\begin{definicja}[BPI i zasada ultrafiltru]\label{def:bpi-principles}
\emph{Twierdzenie o ideale pierwszym w algebrze Boole'a}
(ang. \emph{Boolean Prime Ideal Theorem}, w skrócie BPI)
głosi: w każdej nietrywialnej algebrze Boole'a każdy
właściwy filtr można przedłużyć do ultrafiltru.
\emph{Zasada ultrafiltru} (UF) głosi: dla każdego zbioru
$S$ każdy właściwy filtr podzbiorów $S$ można
przedłużyć do ultrafiltru na $\mathcal P(S)$.
Oba zdania formułujemy nad ZF, nie zakładając AC.
\end{definicja}

\begin{twierdzenie}[Równoważność BPI i UF w ZF]
\label{thm:bpi-ultrafilter-equivalence}
W teorii ZF zasada BPI jest równoważna zasadzie UF.
\end{twierdzenie}
\begin{proof}
\emph{BPI $\Rightarrow$ UF.} Jeśli $S=\varnothing$,
nie istnieje właściwy filtr na $\mathcal P(S)$,
więc twierdzenie jest puste.
Dla niepustego $S$ stosujemy BPI
do nietrywialnej algebry $\mathcal P(S)$. Otrzymane przedłużenie
właściwego filtru jest dokładnie ultrafiltrem podzbiorów.

\emph{UF $\Rightarrow$ BPI, etap 1: zwartość kostki.}
Niech $I$ będzie dowolnym zbiorem. Z UF pokażemy, że
$2^I=\{0,1\}^I$ z topologią produktu dyskretnej pary
jest zwarte. Załóżmy, że rodzina $\mathcal C$ domkniętych
podzbiorów $2^I$ ma własność przecięcia skończonego,
czyli każdy skończony przekrój jest niepusty.
Generuje ona właściwy filtr na $2^I$; UF przedłuża go
do ultrafiltru $U$. Dla każdego $i\in I$ dokładnie jeden
z dwóch cylindrów $\{x:x(i)=0\}$ i $\{x:x(i)=1\}$
należy do $U$. Definiujemy $z(i)$ jako tę jedyną wartość,
której cylinder należy do $U$. Zastępowanie tworzy jedną
funkcję $z:I\to\{0,1\}$, bez żadnych wyborów
z nieokreślonych rodzin. Każde podstawowe otoczenie $z$
to skończony przekrój cylindrów należących do $U$,
więc również należy do $U$. Jeśli $C\in\mathcal C$ i
$z\notin C$, otwarte $2^I\setminus C$ zawierałoby takie
otoczenie $z$, a więc i ono należałoby do $U$, choć
$C\in U$: sprzeczność. Stąd $z\in\bigcap\mathcal C$.
To jest kryterium zwartości przez rodzinę domkniętą.

\emph{Etap 2: zapis aksjomatów algebry w kostce.}
Ustalmy nietrywialną algebrę $B$ i właściwy filtr $F$.
Rozważmy kostkę $2^B$, której punktem jest dowolna
funkcja $h:B\to\{0,1\}$ (niekoniecznie homomorfizm).
Dla $b,c\in B$ narzućmy warunki
\[
 h(0_B)=0,\quad h(1_B)=1,\quad
 h(\neg b)=1-h(b),\quad
 h(b\wedge c)=h(b)h(c),\quad h(b)=1\ (b\in F).
\]
Każda równość zależy od skończenie wielu współrzędnych
kostki, więc wyznacza podzbiór domknięty (nawet
otwarto-domknięty). Pokażemy, że cała rodzina tych
warunków ma własność przecięcia skończonego.

\emph{Etap 3: spełnialność skończonej części.}
Ustalmy skończenie wiele warunków i weźmy skończoną
podalgebrę $D\subseteq B$ generowaną przez wszystkie
występujące w nich elementy. Jest skończona, gdyż jej
elementy są kombinacjami boolowskimi skończenie wielu
generatorów; każda taka kombinacja jest sumą pewnych
spośród skończenie wielu elementarnych koniunkcji.
Niech $f_1,\dots,f_m\in F$ będą elementami,
którym nakazano wartość $1$. Ich iloczyn
$b=f_1\wedge\cdots\wedge f_m$ jest niezerowy,
bo $F$ jest właściwy (dla $m=0$ bierzemy $b=1_B$).
W skończonej algebrze $D$ istnieje atom $a\leq b$:
wśród niezerowych elementów $D$ leżących poniżej $b$
wybieramy element minimalny.
Jeżeli $d\in D$, atomowość $a$ oznacza,
że $a\wedge d$ jest równe $0_B$ lub $a$;
w pierwszym przypadku $a\le\neg d$,
a w drugim $a\le d$.
Nie mogą zachodzić oba przypadki, gdyż wtedy
$a\le d\wedge\neg d=0_B$, sprzeczność.
Dla $d\in D$ połóżmy
$h_D(d)=1$ dokładnie wtedy, gdy $a\leq d$.
Z wykazanej alternatywy wynika
$h_D(\neg d)=1-h_D(d)$.
Ponieważ $a\le d\wedge e$
wtedy i tylko wtedy, gdy
$a\le d$ oraz $a\le e$, mamy
$h_D(d\wedge e)=h_D(d)h_D(e)$.
Ponadto $h_D(0_B)=0$, $h_D(1_B)=1$
i $h_D(f_j)=1$, bo $a\le b\le f_j$. Rozszerzamy $h_D$
do punktu $2^B$, przyjmując na współrzędnych poza $D$
dowolną ustaloną wartość, na przykład zero.
Ten punkt spełnia wszystkie \emph{rozważane skończenie}
warunki; nie zakładaliśmy, że rozszerzenie jest
homomorfizmem całej algebry.

\emph{Etap 4: przejście do wszystkich warunków.}
Zwartość $2^B$ daje funkcję $h$ spełniającą wszystkie
równości jednocześnie. Z pierwszych trzech rodzajów
warunków wynika, że $h:B\to\{0,1\}$ jest homomorfizmem
algebr Boole'a (suma wynika z dopełnień i iloczynu).
Zbiór $U=h^{-1}[\{1\}]$ jest właściwym filtrem;
ponieważ $h(\neg b)=1-h(b)$, Lemat
\ref{lem:bpi-ultrafilter-criterion} mówi, że jest
ultrafiltrem. Ostatnie warunki dają $F\subseteq U$.
To kończy dowód UF$\Rightarrow$BPI.
\end{proof}

\begin{przyklad}[Ultrafiltr niegłówny na $\N$]
\label{ex:bpi-free-ultrafilter}
Filtr wszystkich zbiorów współskończonych w $\N$
jest właściwy. BPI przedłuża go do ultrafiltru $U$.
Taki ultrafiltr nie jest postaci
$U_n=\{A\subseteq\N:n\in A\}$: dla każdego $n$
zbiór $\N\setminus\{n\}$ należy do początkowego filtru,
a zatem i do $U$, podczas gdy nie należy do $U_n$.
Dostaliśmy ultrafiltr niegłówny bez podawania
żadnego wzoru wymieniającego jego elementy.
\end{przyklad}

\section{Dwie wersje twierdzenia Tychonoffa}
\label{sec:weak-choice-tychonoff}
\begin{definicja}[Zwartość i topologia produktu]
\label{def:tychonoff-topology}
Przestrzeń topologiczna $X$ jest \emph{zwarta}, jeżeli
z każdego jej pokrycia zbiorami otwartymi można wybrać
skończone podpokrycie. \emph{Hausdorffa} oznacza, że
każde dwa różne punkty mają rozłączne otoczenia otwarte.
Dla rodziny $(X_i)_{i\in I}$ \emph{topologia produktu}
na $X=\prod_{i\in I}X_i$ ma bazę złożoną ze zbiorów
$\bigcap_{i\in J}\pi_i^{-1}[V_i]$, gdzie
$J\subseteq I$ jest skończony, $V_i\subseteq X_i$ jest
otwarty, a $\pi_i(x)=x(i)$ jest projekcją współrzędnej.
Nie zakładamy, że sam produkt jest niepusty.
\end{definicja}

\begin{twierdzenie}[Tychonoff dla zwartych przestrzeni Hausdorffa]
\label{thm:bpi-tychonoff-hausdorff}
W ZF zasada BPI jest równoważna następującemu zdaniu:
\emph{każdy produkt dowolnej rodziny zwartych przestrzeni
Hausdorffa jest zwarty}.
\end{twierdzenie}
\begin{proof}
\emph{BPI $\Rightarrow$ zwartość produktów.}
Na mocy Twierdzenia~\ref{thm:bpi-ultrafilter-equivalence}
możemy korzystać z UF. Niech
$X=\prod_{i\in I}X_i$; jeśli $X=\varnothing$, to jest
zwarty. W przeciwnym razie rozważmy dowolny ultrafiltr
$\mathcal U$ podzbiorów $X$. Jego obraz przez $\pi_i$
to ultrafiltr $\mathcal U_i$ na $X_i$:
$A\subseteq X_i$ należy do $\mathcal U_i$
wtedy i tylko wtedy, gdy $\pi_i^{-1}[A]\in\mathcal U$.
Zwartość $X_i$ zapewnia istnienie punktu skupienia
$\mathcal U_i$, a Hausdorffowość jego jednoznaczność.
Dokładniej, domknięcia wszystkich zbiorów należących
do $\mathcal U_i$ mają własność przecięcia skończonego;
ich przekrój jest niepusty ze zwartości. Jeśli $x_i$
należy do tego przekroju i $V$ jest jego otwartym
otoczeniem, to $V\in\mathcal U_i$: w przeciwnym razie
$X_i\setminus V\in\mathcal U_i$, a jego domknięcie
nie zawierałoby $x_i$. W przestrzeni Hausdorffa dwa
różne punkty nie mogą być granicami jednego filtru,
więc $x_i$ jest jedyny.

Dla każdego $i$ otrzymaliśmy tym samym \emph{jednoznacznie}
określony $x_i$; funkcja $i\mapsto x_i$ istnieje
przez zastępowanie, bez aksjomatu wyboru. Wektor
$x=(x_i)_{i\in I}$ należy do $X$.
Każde podstawowe otoczenie $x$ jest skończonym
przekrojem przeciwobrazów otoczeń $x_i$, a zatem
należy do $\mathcal U$. Ultrafilter zbiega do $x$.
Aby wyprowadzić zwartość, przypuśćmy, że istnieje
pokrycie otwarte $X$ bez skończonego podpokrycia.
Dopełnienia jego elementów tworzą rodzinę domkniętych
zbiorów z własnością przecięcia skończonego; z UF
rozszerzamy generowany przez nie filtr do ultrafiltru.
Otrzymany ultrafiltr zbiega do $x$, lecz jakiś element
pokrycia zawiera $x$ i tym samym należy do ultrafiltru,
razem ze swoim dopełnieniem. Sprzeczność.

\emph{Zwartość produktów $\Rightarrow$ BPI.}
Załóżmy, że każdy produkt zwartych przestrzeni
Hausdorffa jest zwarty. W szczególności dla każdego
zbioru $B$ kostka $2^B$ jest zwarta, gdyż każdy czynnik
$\{0,1\}$ z topologią dyskretną jest skończony,
a więc zwarty i Hausdorffa. Stosujemy dokładnie
konstrukcję warunków zamkniętych z etapów 2--4
dowodu Twierdzenia~\ref{thm:bpi-ultrafilter-equivalence}:
zwartość kostki daje homomorfizm $h:B\to\{0,1\}$
przyjmujący $1$ na zadanym filtrze. Jego przeciwobraz
$\{1\}$ jest żądanym ultrafiltrem.
\end{proof}

\begin{twierdzenie}[Pełne twierdzenie Tychonoffa i AC]
\label{thm:ac-full-tychonoff}
Nad ZF następujące dwa zdania są równoważne:
(i) aksjomat wyboru AC;
(ii) produkt dowolnej rodziny przestrzeni zwartych
(niekoniecznie Hausdorffa) jest zwarty.
\end{twierdzenie}
\begin{proof}
\emph{AC $\Rightarrow$ (ii).}
AC daje lemat Kuratowskiego--Zorna
(Twierdzenie~\ref{thm:choice-equivalence}), a z niego
BPI: zbiór właściwych filtrów zawierających zadany filtr,
uporządkowany przez zawieranie, ma niepusty łańcuchowy
warunek górnego ograniczenia, bo suma niepustego łańcucha
właściwych filtrów jest właściwym filtrem, a pusty
łańcuch jest ograniczony przez zadany filtr. Zatem każdy
filtr podzbiorów można przedłużyć do ultrafiltru.
Niech $X=\prod_i X_i$ będzie produktem zwartych
przestrzeni. Dla ultrafiltru $\mathcal U$ na $X$ jego
obrazy $\mathcal U_i$ mają w każdym czynniku co najmniej
jeden punkt graniczny. AC pozwala \emph{jednocześnie}
wybrać punkty $x_i$ graniczne dla wszystkich $i$.
Wtedy $x=(x_i)_i$ jest granicą $\mathcal U$ w produkcie,
bo każde jego podstawowe otoczenie zależy od skończenie
wielu współrzędnych. Argument z pokryciem i UF z
poprzedniego dowodu zapewnia zwartość $X$.

\emph{(ii) $\Rightarrow$ AC.}
Niech $(A_i)_{i\in I}$ będzie dowolną rodziną
niepustych zbiorów. Dla każdego $i$ dołączmy jeden
nowy, nazwany punkt $\infty_i\notin A_i$ i połóżmy
$X_i=A_i\sqcup\{\infty_i\}$. Nadajmy $X_i$
topologię $\tau_i=\{\varnothing,A_i,\{\infty_i\},X_i\}$.
Każdy $X_i$ jest zwarty: ma tylko cztery zbiory otwarte;
a ponadto $A_i$ jest domknięto-otwarty. Produkt
$X=\prod_iX_i$ jest niepusty dzięki funkcji
$i\mapsto\infty_i$, nawet jeśli nie mamy wyboru
z rodziny $(A_i)$. Założenie (ii) daje zwartość $X$.
Dla każdego $i$ zbiór
$C_i=\{x\in X:x(i)\in A_i\}$ jest domknięty.
Każdy skończony przekrój zbiorów $C_i$ jest niepusty:
dla skończenie wielu indeksów wybieramy po jednym elemencie
z odpowiednich $A_i$, a na pozostałych współrzędnych
stawiamy $\infty_i$. Zwartość daje
$x\in\bigcap_{i\in I}C_i$. Wtedy $x(i)\in A_i$
dla każdego $i$, czyli $x$ jest funkcją wyboru.
\end{proof}

\begin{uwaga}[Dlaczego Hausdorffowość ma znaczenie?]
Dowód przy BPI wykorzystał \emph{jedyność} granicy
ultrafiltru w przestrzeni Hausdorffa, dzięki czemu
zebraliśmy wszystkie współrzędne w jedną funkcję
bez AC. W dowodzie dla ogólnych przestrzeni zwartych
granic może być wiele; wybranie jednej dla każdego
czynnika może wymagać pełnego AC.
Przestrzenie $X_i$ z ostatniego dowodu są dla
$|A_i|>1$ niehausdorffowskie: dowolne dwie różne
wartości z $A_i$ mają te same nietrywialne otoczenia.
\end{uwaga}

\section{Lemat Teichmüllera--Tukeya i dokładna siła wyboru}
\label{sec:weak-choice-tukey}
\begin{definicja}[Rodzina skończonego charakteru]
\label{def:tukey-finite-character}
Rodzina $\mathcal F\subseteq\mathcal P(E)$ ma
\emph{skończony charakter}, jeżeli dla każdego
$X\subseteq E$ zachodzi równoważność
\[
X\in\mathcal F\quad\Longleftrightarrow\quad
\text{każdy skończony podzbiór }Y\subseteq X\text{ należy do }\mathcal F.
\]
\emph{Lemat Teichmüllera--Tukeya} orzeka, że każda
niepusta rodzina skończonego charakteru ma element
maksymalny względem zawierania.
\end{definicja}

\begin{twierdzenie}[Równoważność lematu Tukeya i AC]
\label{thm:tukey-equivalent-ac}
W teorii ZF lemat Teichmüllera--Tukeya jest
równoważny pełnemu aksjomatowi wyboru.
\end{twierdzenie}
\begin{proof}
\emph{AC $\Rightarrow$ Tukey.}
Niech $\mathcal F\subseteq\mathcal P(E)$ będzie niepustą
rodziną skończonego charakteru, uporządkowaną przez
zawieranie. Dla każdego łańcucha $\mathcal C\subseteq\mathcal F$
jego suma $M=\bigcup\mathcal C$ należy do $\mathcal F$:
każdy skończony $Y\subseteq M$ jest zawarty w pewnym
$X\in\mathcal C$, bo finitarnie wiele jego elementów
pochodzi z finitarnie wielu elementów łańcucha;
wystarczy wybrać największy z tych skończenie wielu
zbiorów łańcucha. Stąd $Y\in\mathcal F$, a skończony
charakter daje $M\in\mathcal F$. Dla pustego łańcucha
suma jest pusta i należy do $\mathcal F$ z jej niepustości
i skończonego charakteru. Z lematu Kuratowskiego--Zorna
(Twierdzenie~\ref{thm:choice-zorn}) istnieje element
maksymalny rodziny.

\emph{Tukey $\Rightarrow$ AC.}
Niech $(A_i)_{i\in I}$ będzie rodziną niepustych zbiorów.
Niech $E=\bigcup_{i\in I}(\{i\}\times A_i)$ i niech
$\mathcal F$ składa się ze wszystkich podzbiorów $F\subseteq E$,
które są grafami funkcji częściowych $f:J\to\bigcup_iA_i$
o dziedzinie $J\subseteq I$, z $f(i)\in A_i$ dla $i\in J$.
Ta rodzina ma skończony charakter: jeżeli relacja
$F\subseteq E$ nie jest grafem funkcji, zawiera dwie
różne pary o tym samym pierwszym współrzędnym;
naruszenie własności jest więc widoczne już w
dwuelementowym podzbiorze. Jeśli każdy skończony
fragment jest grafem funkcji, to jest nim całe $F$.
Z Tukeya istnieje graf maksymalny $F_\ast\in\mathcal F$.
Gdyby jakiś indeks $i\in I$ nie należał do
$\operatorname{dom}F_\ast$, z niepustości $A_i$ moglibyśmy
wybrać jeden element $a\in A_i$ i dołączyć $(i,a)$,
otrzymując większą funkcję częściową. Sprzeczność.
Zatem $F_\ast$ jest grafem funkcji wyboru o dziedzinie $I$.
\end{proof}

\begin{uwaga}[BPI nie jest pełnym wyborem]
\label{rem:bpi-vs-ac}
Mamy w ZF łańcuch implikacji
\[
\mathrm{AC}\Longrightarrow \mathrm{BPI}
\Longleftrightarrow\mathrm{UF}
\Longleftrightarrow\text{Tychonoff dla zwartych Hausdorffa}.
\]
Lemat Tukeya jest \emph{równoważny AC}, a więc
znacznie silniejszy od BPI. Twierdzenie, że BPI nie
implikuje AC nad ZF, jest wynikiem modelowym
(teoria ZF+BPI+$\neg$AC jest względnie niesprzeczna
przy odpowiednich założeniach) i nie wynika z samego
wyżej dowiedzionego łańcucha implikacji.
\end{uwaga}

\begin{uwaga}[Zakres zasady wyboru w dowodzie Tychonoffa]
\label{rem:tychonoff-choice-distinction}
Zwartość pojedynczej przestrzeni nie wymaga wyboru punktu
z każdego czynnika produktu. W kierunku
\emph{BPI$\Rightarrow$Tychonoff dla Hausdorffa} punkt graniczny
ultrafiltru jest \emph{jednoznaczny}, co pozwala utworzyć
cały wektor przez zastępowanie. W ogólnym przypadku
zwartym, ale nie-Hausdorffa, zbiór punktów granicznych
może mieć wiele elementów: do jednoczesnego wyboru
potrzeba AC. To jest matematyczna przyczyna różnicy
między dwiema wersjami twierdzenia, a nie tylko
subtelność użytej techniki dowodu.
\end{uwaga}